\section{Combinations with repetition}

Sometimes we choose \(k\) objects from \(n\) types, and the same type may be chosen several times. The order of the chosen types is ignored. The outcome is then not a sequence but a vector of counts.

\begin{definitionbox}
A \textbf{combination with repetition} of size \(k\) from \(n\) types is a choice of non-negative integers
\[
(x_1,x_2,\ldots,x_n)
\]
such that
\[
x_1+x_2+\cdots+x_n=k.
\]
Here \(x_i\) records how many objects of type \(i\) are chosen.
\end{definitionbox}

The number of such choices is
\[
\binom{n+k-1}{k}=\binom{n+k-1}{n-1}.
\]

This is often called the stars and bars formula. The idea is to represent the \(k\) chosen objects as stars and use \(n-1\) bars to separate the counts of the \(n\) types.

\begin{examplebox}
How many ways are there to choose \(5\) donuts from \(3\) flavors if only the numbers of each flavor matter?

An outcome is a count vector
\[
(x_1,x_2,x_3),\qquad x_1+x_2+x_3=5.
\]
The number of outcomes is
\[
\binom{3+5-1}{5}=\binom{7}{5}=\binom{7}{2}.
\]
\end{examplebox}

\begin{remarkbox}
Combinations with repetition do not describe ordered repeated choices. If order matters, the model is a sequence with repetition and the count is \(n^k\). If order is ignored and only counts of types matter, the model is combinations with repetition.
\end{remarkbox}
